\scanpage{26}
un sens les foncteurs ``ombres'' $\unsure{w|\omega}_N$ sont d'autant plus grossiers que $N$ est plus grand. Néanmoins ils sont tous pleinement fidèles (\uline{NB} les foncteurs $\tau_N$ ne sont pas pl. fidèles bien sur, pas $=$ sur la sous-catégorie des ``ombres virtuelles d'ordre $N-1$'', i.e. des modules gradués dont les composantes sont nulles en degrés $\leq N-1$.

Soit $\operatorname{Omb}(k)$ la catégorie des ``ombres virtuelles'' de $k$-modules, i.e des $\deleted{\unsure{C}}S=k\{T\}$-modules gradués [à degrés $\geq0$ ?]. Soit, $\forall N\geq0$, $\operatorname{Omb}_{N\text{-}\mathrm{tr}}(k)$ la sous-catégorie pleine des modules $N$-Négligeables, i.e. tels que $\mathcal M_i=0$ si $i\geq N$. C'est une sous-catégorie épaisse, la catégorie quotient est can. équivalente \deleted{à celle} \inserted{à la ss-cat. pleine de (Omb) formée} des ``\inserted{$N$-}Ombres virtuelles'', i.e. des $\deleted{\unsure{k}}S$-modules gradués tels que $\mathcal M_i=0$ si $i<N$.

Le foncteur passage au quotient s'identifie alors au foncteur ``troncature'' $\tau_N$. On notera que
\[
\tau_N(\mathcal M)\lhook\joinrel\longrightarrow\mathcal M
\]
(inclusion fonctorielle) \inserted{correspondant au fait que} le foncteur \deleted{\unsure{quo}} passage au quotient admet un adjoint à droite (p. fid. comme de juste)
\[
\operatorname{Hom}_{\mathrm{Omb}}(\mathcal N,\tau_N(\mathcal M))\simeq\operatorname{Hom}_{\substack{N\text{-}\mathrm{Omb}\\\text{ou }\mathrm{Omb}}}(\tau_N(\unsure{\mathcal M|\mathcal N}),\tau_N(\mathcal M)).
\]

Appelons ``négligeable'' \unsure{une ombre|toute ombre|les ombres} $\mathcal M$ qui est $\underset{\longrightarrow}{\lim}$ \inserted{filtrée} d'ombres $N$-négligeables (pour $N$ variable), i.e. \unsure{telle|telles} que \inserted{$\forall i\in\mathbb{Z}$,} $\forall x\in\mathcal M_i$, $\exists N\in\mathbb{N}$ tel que $T^{(\alpha)}.x=0$ pour $\alpha\geq N$ i.e. $S^{+(N)}.x=0$. C'est encore une sous-catégorie épaisse, et le localisé est
% archive sheet number: 25

